\section{Implemented nested orbit-address fusion}
\label{sec:recovery-strengthen3-wrapper}
\paragraph{Status: PROVED/EXACT for supplied same-field affine interfaces.}
LOCALBIT1's declared field query $F(L_b(A^iv))$ and ORBITBATCH1's
shared-parent construction admit a concrete nested implementation.
For $E_j(z)=B_jz+c_j$, set
\[
 M_0=I,\quad b_0=0,\qquad
 M_j=M_{j-1}B_j,\quad b_j=M_{j-1}c_j+b_{j-1}.
\]
Then $F(E_1(\cdots E_k(z)))=F(M_kz+b_k)$; both updates use the old
matrix. Alternatively, evaluate the address chain in reverse order
without precomposition. Under dimension halving with a fixed bound on
same-scale stages, this takes $O(D^2)$ field operations and $O(D\log p)$
coordinate scratch. Descriptor storage, matrix acquisition, field
arithmetic bit costs and all live root-evaluator storage are additional
charged resources. There is one root call per uncached requested field
value, not a free root call.

This result applies only before the final Hadamard mask. An intermediate
re-encoding, change of field, new low-degree extension or SAT return
requires its own interface theorem. Arbitrary affine mixing also requires
rechecking individual-degree eligibility. No vMS extractor parameter
identity or SAT compression theorem is inferred from these identities.

The implementation passes 12,112 bit checks across 1,928 declared
restriction/orbit paths. Independent tests cover 680 general affine
chains and 550 probe-recurrence cases. These are finite interface checks,
not all-input SAT coverage. The report is
\nolinkurl{research/agents/proof/2026-10-08_2102_RECOVERY-STRENGTHEN3.tex}.
It also corrects EXTEVAL1's historical bibliography: the CCC 2019 paper
\emph{Typically-Correct Derandomization for Small Time and Space} is by
William M. Hoza. Its Theorem 5 and Appendix A supply small-space extractor
evaluation on a supplied source, not cheap acquisition from SAT.
P versus NP remains OPEN.
